\section{Renewed closure audit: coherent bypass and global accessibility}
\label{ca:section}

\noindent\textbf{Verdict: the unrestricted channel target remains OPEN.}
The demand for a full closure does not discharge a mathematical hypothesis.
This section tests whether the all-superposition passage and absolute-bypass
conditions can supply the missing resource bound in Section~\ref{ch:section}.
They do supply a global feasibility restriction.  They do not impose a finite
positive-flux dimension on local Maxwell solutions, even when bypass is
certified on the entire transverse plane.  The distinction is proved below,
with an explicit construction and its failure of radiative accessibility.

\subsection{A global power-budget restriction}

\begin{proposition}[PROVED: a necessary threshold condition]
\label{ca:powerbudget}
Consider the compact-scatterer vacuum-exterior class of
Assumptions~\ref{pr:class}--\ref{pr:scattering}, with a legitimate integrable
flux on the full plane $\Pi=\{z=0\}$.  Use the usual radiation normalization:
the angular incoming and outgoing far-field amplitudes $a_\infty,b_\infty$
have squared norms equal to their powers, with
$\|b_\infty\|^2\leq\|a_\infty\|^2=\|a\|^2$.
Then
\begin{equation}
 -\|a\|^2\leq P_\Pi(a)\leq\|a\|^2,
 \qquad
 |P_D(a)|\leq(1+\delta)\|a\|^2
 \label{ca:powerbudgeteq}
\end{equation}
whenever the absolute bypass is at most $\delta\|a\|^2$.
Consequently a nonzero counted subspace requires
\begin{equation}
 \eta\tau\leq1+\delta.
 \label{ca:thresholdfeasibility}
\end{equation}
In particular, its dimension is zero if $\eta\tau>1+\delta$.
\end{proposition}
\begin{proof}
Apply the real Poynting conservation law in a large upper half-ball.
Its flat face has outward normal $-\hat z$.  On passing to infinity,
the hemispherical outward flux is the outgoing upper-hemisphere power
minus the incoming upper-hemisphere power.  Thus
\begin{equation}
 P_\Pi(a)=\|\mathbf1_{\Omega_+}b_\infty\|^2
          -\|\mathbf1_{\Omega_+}a_\infty\|^2.
 \label{ca:hemisphere}
\end{equation}
Here the far-field assertion can be checked directly: the magnetic
amplitudes of outgoing and incoming waves are respectively
$Z_0^{-1}\hat r\times E_{\rm out}$ and
$-Z_0^{-1}\hat r\times E_{\rm in}$.  The incoming--outgoing cross terms
in the real radial Poynting flux cancel, leaving their squared-power
difference.  The definition of the radiation power amplitudes includes
this integrated far-field limit.  Integrability of the plane flux permits
the disk integral to converge to $P_\Pi$.
Each term on the right of \eqref{ca:hemisphere} is nonnegative, and no
larger than $\|a\|^2$, proving the first bound.  Also
$|P_D-P_\Pi|\leq\int_{\Pi\setminus D}|S_z|\leq\delta\|a\|^2$.
Finally the counted inequalities give
$\eta\tau\|a\|^2\leq P_D(a)\leq(1+\delta)\|a\|^2$.
\end{proof}

\begin{remark}[PROVED: Loewner bounds do not count dimension]
\label{ca:loewner}
On any finite counted input space let $Q_D$ be its Hermitian section-flux
matrix.  Polarization and Proposition~\ref{ca:powerbudget} give
\begin{equation}
 \eta T^*T\preceq Q_D\preceq(1+\delta)I,
 \qquad \tau I\preceq T^*T\preceq I.
 \label{ca:loewnereq}
\end{equation}
For $0<\eta\leq1$, the choice $T=I_N$, $Q_D=I_N$, with zero signed and
absolute bypass, satisfies these algebraic inequalities for every $N$.
This is not a physical construction.  It establishes precisely that the
coherent Loewner inequalities and global power budget alone supply no rank
bound.  A spatial Maxwell-to-input accessibility restriction is necessary.
\end{remark}

\subsection{Arbitrarily many local fields with whole-plane coherent bypass control}

The next theorem strengthens the local evanescent example of
Proposition~\ref{cx:evanescent}: it certifies absolute bypass on all of
$\Pi\setminus D$, rather than only vanishing flux on the sides of a
finite collar.  Its normalization is total \emph{local plane flux}, not
power supplied by asymptotic incoming channels.

\begin{theorem}[PROVED: coherent local passage with arbitrarily small bypass]
\label{ca:localpackets}
Fix $k_0>0$, $L>0$, $h>0$, an integer $N\geq1$, and $0<\delta<1$.
Let $D=(-L/2,L/2)^2$ and
$\Omega_h=\mathbb R^2\times(-h,h)$.  There are $N$ linearly independent,
smooth, source-free vacuum Maxwell fields on $\Omega_h$, with finite
weighted field norm there, such that every coefficient vector
$a\in\mathbb C^N$ has
\begin{equation}
 S_z(a;x,y,0)\geq0,\quad
 P_\Pi(a)=\|a\|^2,\quad
 P_D(a)\geq(1-\delta)\|a\|^2,\quad
 \int_{\Pi\setminus D}|S_z(a)|\leq\delta\|a\|^2.
 \label{ca:packetflux}
\end{equation}
All tangential Fourier supports can be chosen above any prescribed
transverse wavenumber.  Thus local source-free Maxwell equations and
all-superposition passage and whole-plane bypass alone do not bound the
dimension, even at fixed section, frequency, and homogeneous permittivity.
\end{theorem}

\begin{proof}
Use the unitary Fourier transform on each transverse coordinate.  For
$b>0$, define
\begin{equation}
 \widehat\phi_b(\xi)=(2b)^{-1/2}\mathbf1_{[-b,b]}(\xi),\qquad
 \phi_b(x)=\frac{\sin(bx)}{\sqrt{\pi b}\,x},\qquad
 g_b(x,y)=\phi_b(x)\phi_b(y),
 \label{ca:sinc}
\end{equation}
with the continuous value at zero.  Plancherel gives
$\|\phi_b\|_2=\|g_b\|_2=1$.  Since $|\sin(bx)|\leq1$,
\begin{equation}
 \int_{|x|\geq L/2}|\phi_b(x)|^2\,dx\leq\frac4{\pi bL},
 \qquad
 \int_{\mathbb R^2\setminus D}|g_b|^2\leq\frac8{\pi bL}.
 \label{ca:sincescape}
\end{equation}
Choose $b\geq8N/(\pi L\delta)$, and choose $Q>k_0$ as large as
desired.  For $j=1,\ldots,N$, put
\begin{equation}
 K_j=Q+3bj,\qquad f_j(x,y)=e^{iK_jx}g_b(x,y),\qquad
 f_a=\sum_{j=1}^Na_jf_j.
 \label{ca:packets}
\end{equation}
Their Fourier supports are disjoint rectangles
$[K_j-b,K_j+b]\times[-b,b]$.  Thus the $f_j$ are orthonormal and
$\|f_a\|_2^2=\|a\|^2$.  Pointwise Cauchy--Schwarz gives
$|f_a|^2\leq N|g_b|^2\|a\|^2$, hence
\begin{equation}
 \int_{\mathbb R^2\setminus D}|f_a|^2\leq\delta\|a\|^2.
 \label{ca:packettail}
\end{equation}
Every support lies in $\xi_x\geq Q+2b>k_0$, entirely in the
evanescent sector, but its union is bounded for each finite $N,Q,b$.

Use the normalized fields $e=\sqrt{\varepsilon_0}E$,
$h=\sqrt{\mu_0}H$ of \eqref{ch:normalization}, and set
$C=\sqrt{2/c_0}$.  Prescribe the tangential data at $z=0$ by
\begin{equation}
 e_t=(0,Cf_a),\qquad h_t=(-Cf_a,0),\qquad
 e_z=\frac{iC}{k_0}\partial_yf_a,\qquad
 h_z=-\frac{iC}{k_0}\partial_xf_a.
 \label{ca:cauchydata}
\end{equation}
The normal components are determined by the two normal curl equations.
For completeness, the full extension is constructed, not assumed.  Put
$J(v_x,v_y)=(-v_y,v_x)$ and evolve the tangential components by
\begin{equation}
 \begin{aligned}
 e_z&=\frac{i}{k_0}(\partial_xh_y-\partial_yh_x),&
 \partial_z e_t&=\nabla_te_z-ik_0J h_t,\\
 h_z&=-\frac{i}{k_0}(\partial_xe_y-\partial_ye_x),&
 \partial_z h_t&=\nabla_th_z+ik_0J e_t.
 \end{aligned}
 \label{ca:cauchyevolution}
\end{equation}
After transverse Fourier transformation this is a four-component constant
coefficient ordinary differential equation for each $\xi$.  Its matrix
entries are polynomials of degree at most two in $\xi$, with fixed
coefficients depending on $k_0$.  On the bounded Fourier support the
matrix exponential and every fixed number of its derivatives are bounded
for $|z|\leq h$.  Inverse Fourier transformation therefore supplies smooth
fields, with every transverse derivative and every $z$ derivative in
$L^2(\Omega_h)$.  In particular their weighted field norm is finite.
The normal definitions and the tangential evolution are exactly all six
curl equations, so
$\operatorname{curl}e=ik_0h$ and
$\operatorname{curl}h=-ik_0e$.  This also proves the divergence conditions.

At the measurement plane,
$S_z=(c_0/2)\operatorname{Re}(e\times\overline h)\cdot\hat z
=(c_0C^2/2)|f_a|^2=|f_a|^2$.
Orthogonality and \eqref{ca:packettail} give all four assertions in
\eqref{ca:packetflux}.  They hold for every vector, not merely for the
basis fields or a sample of random inputs.  Nonzero $f_a$ also proves
linear independence of the fields.
\end{proof}

\begin{proposition}[PROVED: the flux normalization permits unbounded field resource]
\label{ca:packetresource}
For the fields in Theorem~\ref{ca:localpackets}, set $q_{\min}=Q+2b$.
Their local slab resource satisfies
\begin{equation}
 \int_{\Omega_h}(|e_a|^2+|h_a|^2)\,dx\,dy\,dz
 \geq\frac{4h}{c_0k_0^2}q_{\min}^2\|a\|^2.
 \label{ca:packetresourceeq}
\end{equation}
Thus the same coherent flux and bypass bounds coexist with arbitrarily
large resource at fixed $N,D,h,k_0,\delta$ in this local class.
\end{proposition}
\begin{proof}
The vacuum Maxwell equations imply $(\Delta+k_0^2)h_z=0$.
For each Fourier frequency in the support, put
$\kappa(\xi)=\sqrt{|\xi|^2-k_0^2}>0$.
The component has the form
$u(z)=u(0)\cosh(\kappa z)+u'(0)\sinh(\kappa z)/\kappa$.
When $|u|^2$ is integrated on $(-h,h)$ its cross term is odd, and both
remaining terms are nonnegative.  Hence
$\int_{-h}^h|u(z)|^2dz\geq2h|u(0)|^2$.
Plancherel and \eqref{ca:cauchydata} yield
\begin{equation}
 \int_{\Omega_h}|h_z|^2\geq
 \frac{2hC^2}{k_0^2}\|\partial_xf_a\|_2^2
 \geq\frac{4h}{c_0k_0^2}q_{\min}^2\|a\|^2,
 \label{ca:packetresproof}
\end{equation}
which is bounded above by the total resource.
\end{proof}

\subsection{Why the local theorem is not the requested counterexample}

\begin{proposition}[PROVED: failure of a vacuum outgoing completion]
\label{ca:outgoingfailure}
No nonzero field of Theorem~\ref{ca:localpackets} can be continued into an
entire homogeneous vacuum upper half-space with the outgoing condition
that every evanescent Fourier component decays as $z\to+\infty$.
\end{proposition}
\begin{proof}
The transverse Fourier support is preserved by vacuum evolution and
contains only $|\xi|>k_0$, with a positive distance from the light circle.
An outgoing continuation therefore has exponentially decaying $L^2$
electric and magnetic slices as $z\to+\infty$, implying
$P_\Pi(z)\to0$ by Cauchy--Schwarz.
On the other hand, the full-plane integrated flux of a source-free vacuum
solution with these $L^2$ slices is independent of $z$.  One way to justify
this is to integrate the divergence-free real Poynting vector over a
transverse cutoff of radius $R$.  Its lateral error is bounded by
$C/R$ times the $L^1$ norm of the transverse Poynting field integrated
over the finite $z$ interval.  That norm is finite by Cauchy--Schwarz,
so the error tends to zero.  Conservation would give
$P_\Pi(z)=P_\Pi(0)=\|a\|^2>0$, a contradiction.
\end{proof}

The example therefore cannot acquire a physical incoming-power
normalization by simply declaring its coefficient space to be an incident
channel space.  Formally appending $T=I_N$ would satisfy throughput and,
for $\eta\leq1-\delta$, passage, but it would not construct the global
field-to-port map required in Assumption~\ref{pr:scattering}.

\begin{proposition}[PROVED: no exact realization of these entire slab data by a compact exterior device]
\label{ca:compactcompletion}
Suppose a proposed dielectric perturbation is compactly supported outside
$\Omega_h$ and the exterior convention consists of propagating incoming
radiation and outgoing scattered radiation.  It cannot have a nonzero
field of Theorem~\ref{ca:localpackets} as its exact total field on all of
$\Omega_h$.
\end{proposition}
\begin{proof}
The Fourier Cauchy construction defines a homogeneous Maxwell solution
$F^0=(e^0,h^0)$ on all of $\mathbb R^3$: for every finite $z$ its bounded
Fourier support makes the matrix exponential well defined.  Let $F$ be
the proposed scattering field.  The difference $F-F^0$ is a homogeneous
vacuum Maxwell field in the unbounded connected component of the
complement of the material support, and is zero on the entire open slab.
Each component satisfies the constant-coefficient Helmholtz equation,
so analytic continuation implies equality throughout that connected
component, in particular everywhere outside a sufficiently large ball.

Choose $z_0$ above that ball.  The entire half-space $z>z_0$ is vacuum,
and $F=F^0$ there.  Its transverse spectrum is the same bounded set
$|\xi|>k_0$ as before.  In a homogeneous exterior half-space the stated
radiation convention permits incoming propagating components, outgoing
propagating components, and evanescent components decaying toward
infinity; it excludes growing evanescent components.  This follows as
well by transverse Fourier transformation of the outgoing vacuum Green
function and of the incident propagating angular spectrum: the
longitudinal root for $|\xi|>k_0$ has positive imaginary part.
There are no propagating components in the present support.  The field
therefore decays toward $+\infty$ in $L^2$ slices, contradicting its
conserved flux $\|a\|^2>0$, exactly as in
Proposition~\ref{ca:outgoingfailure}.
\end{proof}

Thus even the exact local packet data are not a ready-made physical
counterexample.  A resonant construction would need a quantitatively
controlled \emph{approximation} to suitable concentrated data, or
different accessible fields, with input power and every coherent error
tracked.  Proposition~\ref{ca:compactcompletion} does not exclude that
possibility, nor does it exclude arbitrary external resonant matching.

\subsection{Matching theory and the decisive missing statement}

The primary reflectionless-scattering analysis of
\href{https://doi.org/10.1515/nanoph-2020-0403}{Stone, Sweeney, Hsu,
Wisal, and Wang (2021)} gives a Maxwell formulation of reflection zeros
and structural tuning to a real-frequency reflectionless state.  Its
multichannel state is a selected coherent input vector.  Such a state
does not assert an $N$-dimensional subspace with uniformly high
transmission at the same frequency, nor does it certify passage and
absolute bypass through a fixed small section.  Applying that framework
therefore still requires the simultaneous synthesis and error control
identified in Remark~\ref{cx:missing-bridge}.

\begin{remark}[OPEN: smallest remaining global bridge]
\label{ca:decisivebridge}
The renewed audit reduces the choice to two substantive global tasks.
A proof must restrict the radiatively accessible trace subspaces of all
admitted bounded-index devices strongly enough to give the uniform
finite-rank estimate of Lemma~\ref{ch:ranklemma}, despite arbitrary
external support and feedback.  A disproof must realize an unbounded
sequence of positive-flux trace subspaces with power-normalized incident
and delivered channels and controlled absolute bypass.  Theorem
\ref{ca:localpackets} shows that local Maxwell equations and even a
whole-plane coherent bypass certificate cannot replace that global step.
Proposition~\ref{ca:powerbudget} only closes the empty parameter regime
$\eta\tau>1+\delta$; it does not resolve the meaningful positive-throughput
regime.  No full fixed-section vector theorem or physically admissible
many-channel counterexample has been obtained.
\end{remark}

\begin{remark}[Formalization boundary]
\label{ca:leanboundary}
Finite-dimensional Loewner consequences, norm perturbation certificates,
and dimension lemmas can be formalized independently.  Such Lean
statements do not formalize the existence of a Maxwell scattering
realization.  The construction in Theorem~\ref{ca:localpackets} requires
Fourier analysis, the stated Maxwell evolution, and flux integration;
its proof here is a mathematical derivation, not a Lean verification.
Introducing an axiom that supplies the global bridge in
Remark~\ref{ca:decisivebridge} would assume the unresolved content rather
than verify it.  Accordingly no full-OQ14 Lean closure is claimed.
\end{remark}
